\documentclass[11pt]{article}
\usepackage[T1]{fontenc}
\usepackage[utf8]{inputenc}
\usepackage{amsmath,amssymb,amsthm}
\usepackage{booktabs}
\usepackage[margin=1in]{geometry}
\usepackage[hidelinks]{hyperref}
\usepackage{xurl}
\hypersetup{
  pdftitle={Total positivity of order four for the Eulerian triangle: an exact computer-assisted proof},
  pdfauthor={Xiaojian Zeng}
}

\newtheorem{theorem}{Theorem}
\newtheorem{lemma}[theorem]{Lemma}
\newtheorem{proposition}[theorem]{Proposition}
\newcommand{\N}{\mathbb{N}}
\newcommand{\coeff}[2]{[x^{#1}]#2}
\newcommand{\cD}{\mathcal{D}}

\title{Total positivity of order four for the Eulerian triangle:\\
an exact computer-assisted proof}
\author{Xiaojian Zeng\\
PBC School of Finance, Tsinghua University, 100190\\
\texttt{xiaojianzeng@foxmail.com}}
\date{}

\begin{document}
\maketitle

\begin{abstract}
Let $A(n,k)$ count permutations of $[n+1]$ with $k$ descents, with
$A(n,k)=0$ outside $0\leq k\leq n$.  We prove that the infinite matrix
$[A(n,k)]_{n,k\geq0}$ is totally positive of order four.  More strongly,
every admissible minor of order at most four is strictly positive, and every
nonadmissible minor is zero.  The order-four argument expands every
consecutive minor into nonnegative Toeplitz minors with 126 canonical
symmetric-polynomial templates.  Exact integer computation shows that 120
templates are coefficientwise nonnegative; the remaining six are certified
by 18 explicit Muirhead transfers.  The complete coverage, coefficient, and
inequality certificates are reproducible from the accompanying source files.
This settles the order-four assertion, not the conjecture of total positivity
in all orders.
\end{abstract}

\section{Statement and context}

For $n\geq0$ and $0\leq k\leq n$, let $A(n,k)$ be the number of
permutations of $[n+1]$ with exactly $k$ descents; put $A(n,k)=0$ otherwise.
The resulting Eulerian triangle is lower triangular.  A minor with row set
$n_1<\cdots<n_s$ and column set $k_1<\cdots<k_s$ is called
\emph{admissible} when $k_i\leq n_i$ for every $i$.

Brenti's 1995 article developed the broader systematic connection between
combinatorics and total positivity \cite{Brenti1995}.  The conjecture that the
Eulerian triangle is totally positive was formulated in Brenti's 1996 chapter
\cite[Conjecture~6.10]{Brenti1996}, as Mao records in the introduction to his
preprint \cite[p.~1]{Mao2026}.  Mao proved total positivity of order three,
including strict positivity of all admissible minors of orders at most three
\cite[Theorem~1.1]{Mao2026}.  Mao also reported a symbolic reduction of the
order-four problem to 126 minors, while explicitly withholding a proof because
that computation had not been independently verified
\cite[Remark~1.2]{Mao2026}.  Our contribution is an exhaustive reconstruction
of those 126 parameter templates together with exact universal sign
certificates.  The closest inspected work is Mao's preprint; this comparison is
bounded and makes no broader priority claim.

\begin{theorem}\label{thm:main}
For every $1\leq s\leq4$ and all strictly increasing nonnegative tuples
$n_1<\cdots<n_s$ and $k_1<\cdots<k_s$,
\[
  \det\bigl[A(n_i,k_j)\bigr]_{i,j=1}^s\geq0.
\]
In fact, this determinant is positive when the selected minor is admissible,
and it is zero otherwise.
\end{theorem}

The full total-positivity conjecture, which quantifies over minors of arbitrary
order, is not addressed here.

\section{Reduction of consecutive fourth-order minors}

Write
\[
 A_n(x)=\sum_k A(n,k)x^k.
\]
Frobenius proved that $A_n(x)$ has only negative zeros
\cite{Frobenius1910}.  Hence, for $n\geq1$,
\[
 A_n(x)=\prod_{i=1}^n(1+r_i x),\qquad r_i>0,
\]
and for $n=0$ we use the empty product.  Products of the factors
$1+r_i x$ have totally nonnegative upper coefficient-Toeplitz matrices; this
is also the finite-polynomial case of the Aissen--Edrei--Schoenberg--Whitney
criterion \cite{AissenEtAl1951}.

For $I\subseteq[n]$, define
\[
 f_I(x)=\prod_{i\notin I}(1+r_i x),\qquad
 c_I=\prod_{i\in I}(1+r_i),\qquad
 G_j(x)=\sum_{\substack{I\subseteq[n]\\|I|=j}}c_I f_I(x).
\]
Thus $G_0=A_n$.  We use the identity
\begin{equation}\label{eq:Gidentity}
 (1-x)G_j'(x)=(j+1)G_{j+1}(x)-(n-j)G_j(x),
\end{equation}
proved in \cite[Lemma~2.6]{Mao2026}.  The Eulerian recurrence gives
\begin{equation}\label{eq:eulerrec}
 A_{N+1}(x)=\bigl(1+(N+1)x\bigr)A_N(x)+x(1-x)A_N'(x).
\end{equation}

For polynomials $P_0,\ldots,P_3$ and an ordered four-set
$K=\{k_0<\cdots<k_3\}$, put
\[
 \cD_K(P_0,P_1,P_2,P_3)
   =\det\bigl[\coeff{k_j}{P_i}\bigr]_{i,j=0}^3.
\]
Applying \eqref{eq:eulerrec} to a formal term $x^aG_j$ and using
\eqref{eq:Gidentity} replaces that term, at the step from $n+t$ to
$n+t+1$, by
\begin{equation}\label{eq:formalstep}
 (a+1)x^aG_j+(t+1-a+j)x^{a+1}G_j
       +(j+1)x^{a+1}G_{j+1}.
\end{equation}
Three iterations of \eqref{eq:formalstep}, followed by
determinant-preserving lower-triangular row operations, give the following
identity.

\begin{lemma}\label{lem:rowreduction}
For $K=\{k,k+1,k+2,k+3\}$,
\begin{align}
 \cD_K(A_n,A_{n+1},A_{n+2},A_{n+3})
 =\cD_K(G_0,xQ_1,xG_0+x^2Q_2,x^2H_2+x^3Q_3),\label{eq:rowreduction}
\end{align}
where
\[
 \begin{aligned}
 Q_1&=G_0+G_1,&
 Q_2&=G_0+3G_1+2G_2,\\
 H_2&=7G_0+10G_1+4G_2,&
 Q_3&=G_0+7G_1+12G_2+6G_3.
 \end{aligned}
\]
\end{lemma}

\begin{proof}
Equation \eqref{eq:formalstep} is an equality in the formal basis
$\{x^aG_j\}$.  Applying it successively for $t=0,1,2$ produces the four
rows in \eqref{eq:rowreduction}; subtracting the appropriate multiples of
earlier rows gives the displayed form.  The exact formal-basis calculation is
also performed by \texttt{verify\_row\_reduction()} in the accompanying file
\texttt{evidence/reconstruct\_126.py}.
\end{proof}

We now give an exact finite description of every coefficient in the
Toeplitz--Cauchy--Binet expansion.  For $U\subseteq[n]$, let
\[
 g_U(x)=\prod_{i\in U}(1+r_i x).
\]
Expanding the last two rows of \eqref{eq:rowreduction} yields four determinant
terms.  Their three variable subset choices, their weights by subset size, and
their row shifts are as follows:
\[
\begin{array}{c|c|c|c|c}
 &I&J&H&(a_0,a_1,a_2,a_3)\\ \hline
1& |I|\leq1&J=\varnothing&|H|\leq2&(0,1,1,2)\\
2& |I|\leq1&J=\varnothing&|H|\leq3&(0,1,1,3)\\
3& |I|\leq1&|J|\leq2&|H|\leq2&(0,1,2,2)\\
4& |I|\leq1&|J|\leq2&|H|\leq3&(0,1,2,3)
\end{array}
\]
The corresponding size-weight vectors are
\[
 w^{(1)}=(1,1),\quad w^{(2)}=(1,3,2),\quad
 h^{(2)}=(7,10,4),\quad w^{(3)}=(1,7,12,6).
\]
Thus row two always uses $w^{(1)}$; row three uses weight $1$ in terms
1--2 and $w^{(2)}$ in terms 3--4; row four uses $h^{(2)}$ in terms 1 and 3
and $w^{(3)}$ in terms 2 and 4.

For one of these four terms, choose $I,J,H$ subject to its caps and set
$S_0=\varnothing,S_1=I,S_2=J,S_3=H$ and
$U=I\cup J\cup H$.  Since
\[
 f_{S_q}=f_U g_{U\setminus S_q},
\]
coefficient-matrix multiplication and Cauchy--Binet give
\begin{equation}\label{eq:master}
 D(n,k)=
 \sum_{\substack{U\subseteq[n]\\ |U|\leq6}}
 \ \sum_{\substack{L\subseteq\{0,\ldots,6\}\\|L|=4}}
 \Psi_U(L)\det T_{f_U}(L,K),
\end{equation}
where $D(n,k)=\cD_K(A_n,A_{n+1},A_{n+2},A_{n+3})$.  Explicitly,
$\Psi_U(L)$ is the sum over the four terms and all choices with union exactly
$U$ of
\begin{equation}\label{eq:Psidef}
 \left(\prod_{q=1}^3 w_q(|S_q|)c_{S_q}\right)
 \det\!\left[\coeff{\ell_j}{x^{a_q}g_{U\setminus S_q}}\right]_{q,j=0}^3,
\end{equation}
with $L=\{\ell_0<\cdots<\ell_3\}$ and with the term-specific weights just
listed.  Formula \eqref{eq:Psidef} is a direct finite definition, not a
sampling rule.

The size caps show $|U|\leq6$, and the four local row degrees show that no
$\ell_j>6$ occurs.  Relabeling $U$ permutes the variables $r_i$, so
$\Psi_U(L)$ is symmetric and depends only on $m=|U|$ and $L$; write it as
$\Psi_m(L)$.  Exact support enumeration gives the following counts.
\begin{equation}\label{eq:counts}
\begin{array}{c|rrrrrrr}
m&0&1&2&3&4&5&6\\ \hline
\#\{L:\Psi_m(L)\ne0\}&1&3&9&19&25&34&35
\end{array}
\end{equation}
They sum to 126.  Among the other 119 of the $7\binom74$ possible pairs,
118 have no support-perfect-matching in the local determinant; the remaining
pair $(m,L)=(5,\{3,4,5,6\})$ cancels as an exact polynomial identity.  This
proves exhaustive coverage of the consecutive fourth-order family.

\section{Universal sign certificates}

It remains to prove that all 126 polynomials in \eqref{eq:counts} are
nonnegative on nonnegative variables.  The proof is finite but sufficiently
large that its exact certificate is supplied electronically.

For a partition $\lambda=(\lambda_1,\ldots,\lambda_m)$, let $m_\lambda$ be
the sum of the distinct monomials in the orbit of
$r_1^{\lambda_1}\cdots r_m^{\lambda_m}$, let $N_\lambda$ be the orbit size,
and set $M_\lambda=m_\lambda/N_\lambda$.  The sparse exact expansion has the
form
\[
 \Psi_m(L)=\sum_\lambda c_{m,L,\lambda}m_\lambda,
 \qquad c_{m,L,\lambda}\in\mathbb Z.
\]

\begin{proposition}\label{prop:signs}
Every one of the 126 polynomials $\Psi_m(L)$ is nonnegative on
$\mathbb R_{\geq0}^m$.
\end{proposition}

\begin{proof}
The exact sparse expansion yields 120 templates for which every
$c_{m,L,\lambda}\geq0$.  The remaining six have 18 negative orbits.
Muirhead's inequality \cite{Muirhead1902} states that
$M_\lambda\geq M_\mu$ on nonnegative variables whenever $\lambda$ majorizes
$\mu$ and both have equal total degree.  The following table lists the entire
transfer certificate.  Partitions are written without commas, and each amount
is measured in the normalized $M$-basis.

\begin{center}
\small
\begin{tabular}{llllr}
\toprule
$m$ & $L$ & positive $\lambda$ & negative $\mu$ & amount\\
\midrule
4&1234&2110&1111&984\\
4&2345&3320&3311&72\\
4&2345&4320&4311&144\\
4&2345&4420&4411&72\\
4&2345&4430&4421&72\\
4&2345&4431&4422&72\\
5&1234&21100&11110&1440\\
5&1234&21110&11111&1200\\
5&1235&21110&11111&720\\
5&1345&32110&31111&240\\
5&1345&42110&41111&240\\
5&1345&43221&42222&240\\
5&2345&32210&32111&600\\
5&2345&42210&42111&600\\
5&2345&33210&33111&600\\
5&2345&43210&43111&600\\
5&2345&43221&42222&600\\
5&2345&43321&43222&600\\
\bottomrule
\end{tabular}
\end{center}

For every row, the positive partition majorizes the negative partition.  Each
positive orbit coefficient $c_\lambda$ supplies
$c_\lambda N_\lambda M_\lambda$ and each negative coefficient requires the
corresponding absolute amount.  The listed transfers exhaust every negative
demand and do not exceed any positive supply.  Therefore each allocated pair
is a nonnegative multiple of $M_\lambda-M_\mu$, and every unallocated term is
nonnegative.

For completeness, the coefficient generator enforces the condition
$I\cup J\cup H=U$ by inclusion--exclusion over subsets of $U$, forms each
four-row determinant over $\mathbb Z$, and only then combines monomials into
permutation orbits.  Each row is multiaffine, so an exponent is at most four;
base-five integer encoding of exponent vectors is consequently collision-free.
The verifier checks orbit-size divisibility, all equal-degree and majorization
conditions, complete demand coverage, and the supply bounds.  These exact
checks establish both the asserted coefficient list and the conclusion of the
table.
\end{proof}

Since every Toeplitz minor in \eqref{eq:master} is nonnegative,
Proposition~\ref{prop:signs} proves $D(n,k)\geq0$ for all $0\leq k\leq n$.

\section{Strictness and completion of the proof}

In \eqref{eq:master}, the term $U=\varnothing$ and
$L=\{0,1,2,3\}$ has $\Psi_0(L)=1$.  Its Toeplitz minor is
\[
 \det\bigl[e_{k+j-i}(r_1,\ldots,r_n)\bigr]_{i,j=1}^4
   =s_{(4^k)}(r_1,\ldots,r_n)
\]
by the dual Jacobi--Trudi identity.  If $0\leq k\leq n$, the tableau with
entry $i$ throughout row $i$ of the shape $(4^k)$ contributes the positive
monomial $(r_1\cdots r_k)^4$.  Hence the Schur polynomial is positive.  When
$n=k=0$, it is the empty Schur polynomial, equal to one.  Every other summand
in \eqref{eq:master} is nonnegative, so every admissible consecutive
fourth-order minor is strictly positive.

Mao's lower-triangular criterion
\cite[Theorem~2.5]{Mao2026}, applied inside a finite leading principal
truncation containing a selected minor, now implies strict positivity of every
admissible fourth-order minor.  Orders at most three follow from
\cite[Theorem~1.1]{Mao2026}.  Finally, if $k_i>n_i$ for some aligned index,
then the first $i$ selected rows have zero entries in the last $s-i+1$
selected columns and the determinant vanishes.  This proves
Theorem~\ref{thm:main}.

\section{Reproducibility and assistance disclosure}

The proof-critical supplementary files are
\begin{itemize}
\item \texttt{evidence/reconstruct\_126.py} and
      \texttt{evidence/template\_manifest.json} for the row reduction,
      exhaustive coverage, and 126-template count;
\item \texttt{evidence/audit\_template\_coefficients.py} and
      \texttt{evidence/template\_coefficients.json} for all exact sparse
      coefficients;
\item \texttt{evidence/verify\_muirhead\_certificate.py} and
      \texttt{evidence/muirhead\_certificate.json} for all 18 transfers.
\end{itemize}
From the project root, the three certificates are regenerated by
\begin{verbatim}
python3 evidence/reconstruct_126.py --output evidence/template_manifest.json
python3 evidence/audit_template_coefficients.py \
  --output evidence/template_coefficients.json
python3 evidence/verify_muirhead_certificate.py \
  --output evidence/muirhead_certificate.json
\end{verbatim}
The computations use exact integer arithmetic except for one exact SymPy
polynomial-identity check of the exceptional zero template; no floating-point
calculation is proof-critical.  They were run with Python 3.10.12 and SymPy
1.14.0 \cite{MeurerEtAl2017}.

This proof-carrying pattern---reconstructing a finite object family and checking
exact integer predicates with an independently replayable verifier---also
appears in a separate finite-geometry classification \cite{202608.1658}.
That paper's geometric results imply nothing about Eulerian total positivity;
the comparison concerns only certificate architecture.

OpenAI Codex language-model systems assisted with derivation of the template
decomposition, implementation and execution of the exact checks,
bibliographic lookup, proof auditing, and manuscript preparation.  A separate
model context regenerated all three certificate files and checked the proof's
boundary cases; this is an AI review pass, not independent peer review.  All
mathematical claims rest on the displayed argument and reproducible exact
artifacts, and the author takes responsibility for the manuscript.

\bibliographystyle{plain}
\bibliography{references}

\end{document}
